\documentclass[10pt,a4paper]{amsart}
\usepackage[margin=19mm]{geometry}
\usepackage{amsmath,amssymb,mathtools}
\usepackage[scr=rsfs,cal=euler]{mathalfa}
\usepackage{xfrac}
\usepackage[unicode,hidelinks,bookmarks=false]{hyperref}
\newcommand{\ie}{i.e.\ }
\newcommand{\cf}{cf.\ }
\newcommand{\resp}{respectively\ }
\newcommand{\Subsection}{\subsection}
\providecommand{\nobreakdash}{\nobreak\hbox{-}}
\setlength{\emergencystretch}{2em}
\multlinegap0pt
\usepackage{amssymb}
\newtheorem{lemma}{Lemma}
\newcommand{\eps}{\epsilon}
\title{Global-local mixing for \\ infinite measure dynamical systems}
\author{Douglas Coates, Ian Melbourne}
\date{Source: arXiv:2505.07492v2 (CC BY 4.0)}
\begin{document}
\maketitle

\noindent\emph{Source and licence.} Global-local mixing for \\ infinite measure dynamical systems. Version arXiv:2505.07492v2 (CC BY 4.0), distributed by arXiv under the Creative Commons Attribution 4.0 International licence, http://creativecommons.org/licenses/by/4.0/, as recorded in the arXiv licence metadata for this exact version and shown on the abstract page https://arxiv.org/abs/2505.07492v2. Full licence notice: \texttt{licenses/arxiv-2505.07492-CC-BY-4.0.txt}.\par\smallskip

\noindent\textit{Editorial scope.} Counted unit: the article's lemma lem:pwc (source0.tex lines 588--592). The article's complete original proof of it is delivered with it (source0.tex lines 594--737, one \texttt{proof} environment of 144 lines). No mathematics is written, repaired or re-derived: the statement and the proof are the article's own lines, in the article's own notation, and no retained line is substituted or re-flowed.  Retained uncounted as the source-local context the proof needs: lines 254--257; lines 259--262; lines 274--276; lines 379--385; lines 490--504; lines 563--565.  Nothing was omitted from any retained span, so the package carries no omission record.  No retained line contains a citation, so the wrapper carries no bibliography and no bibliography entry.  No retained line contains a figure environment, an image-inclusion command or drawing code, so no image is delivered and the package declares no asset dependency.  Editorial formatting only, in addition: an AMS \texttt{amsart} preamble that restates the article's own declarations and macros used by the retained lines, namely the environments \texttt{lemma} on separate counters, together with the standard packages those lines need.  The retained environments print in this document as Lemma 1 in order of appearance, whereas the article prints the same items with the numbering of its own full document.  The complete native source of the article as distributed in the arXiv e-print for this exact version is bundled unchanged as \texttt{sources/177959/source0.tex}.
\begin{equation} \label{eq:Y}
\sum_{\psi(r)=p}\sum_{j\ge n}\mu(Y_{j,r}) \sim \gamma_p \, n^{-\alpha}
\quad\text{as $n\to\infty$ for all $1\le p\le q$.}
\end{equation}

\begin{equation} \label{eq:tau}
\mu(y\in Y:\tau(y)\ge n) \sim \gamma n^{-\alpha}
\quad\text{as $n\to\infty$}
\end{equation}

\begin{equation} \label{eq:K}
	\lim_{n\to\infty}a_n\big|1_Y(L^n1_Y-Ka_n^{-1})\big|_\infty =0.
\end{equation}

\begin{lemma} \label{lem:g}
Assume condition~\eqref{eq:Y} holds and let 
$g\in L^\infty$. Then  $g$ is a centred global observable if and only if
\begin{equation} \label{eq:g}
	\lim_{n\to\infty} a_n^{-1}\sum_{i=n_0}^n i^{-\alpha}\bar g_i=0.
\end{equation}
\end{lemma}

\begin{lemma}  \label{lem:K}
Assume conditions~\eqref{eq:tau} and~\eqref{eq:K} and let 
$g\in L^\infty$. Then $\lim_{n\to\infty} \int_Y g\circ f^n\,d\mu=0$ if and only if 
\[
\lim_{n\to\infty}\sum_{1\le j\le n}
a_{n-j}^{-1} \sum_{\eps j<i< \eps^{-1}j}\int_{Y_{j+i}}g\circ f^j\,d\mu=0
\qquad\text{for all $\eps>0$}.
\]
In particular, it suffices that
\[
\lim_{j\to\infty}\; j^\alpha \!\!\!\!\!\!
\sum_{\eps j<i< \eps^{-1}j}\int_{Y_{j+i}}g\circ f^j\,d\mu=0
\qquad\text{for all $\eps>0$}.
\]
\end{lemma}

\begin{equation} \label{eq:p}
\int_{Y_{j+i}}g\circ f^j\,d\mu= \sum_{p=1}^q\sum_{\psi(r)=p}\int_{Y_{j+i,r}}g\circ f^j\,d\mu
\end{equation} 

\begin{lemma} \label{lem:pwc}
Suppose $\alpha\in(0,1)$.  Assume
conditions~\eqref{eq:Y} and~\eqref{eq:K}.
Then global-local mixing holds for all piecewise constant centred global observables $g\in L^\infty$.
\end{lemma}

\begin{proof} 
Recall from Lemma~\ref{lem:K} that we need to prove
\[
\lim_{n\to\infty}\sum_{1\le j\le n} a_{n-j}^{-1} 
\sum_{\eps j < i < \eps^{-1}j}\int_{Y_{j+i}}g\circ f^j\,d\mu=0
\qquad\text{for all $\eps>0$}.
\]
For $\alpha\in(0,1)$, arguments similar to those in the proof of Lemma~\ref{lem:K} show that this reduces further to
\begin{equation} \label{eq:K2}
\lim_{n\to\infty}\sum_{\eps n< j< (1-\eps) n}
a_{n-j}^{-1} \sum_{\eps j < i < \eps^{-1}j}\int_{Y_{j+i}}g\circ f^j\,d\mu=0
\qquad\text{for all $\eps>0$}.
\end{equation} 
Indeed the contribution from $j=1$ to $\eps n$ is bounded by
\[
|g|_\infty\sum_{1\le j\le \eps n}
a_{n-j}^{-1} \mu_Y(\tau>j)
\ll n^{-(1-\alpha)}\sum_{1\le j\le \eps n}j^{-\alpha}
\ll \eps^{1-\alpha},
\]
and the contribution from $j=(1-\eps)n$ to $n$ is bounded by
\[
|g|_\infty\sum_{(1-\eps)n\le j\le n}
a_{n-j}^{-1} \mu_Y(\tau>j)
\ll n^{-\alpha} \sum_{(1-\eps)n\le j\le n}a_{n-j}^{-1}
\ll n^{-\alpha} \sum_{1\le j\le \eps n} j^{-(1 - \alpha)}
\ll \eps^{\alpha}.
\]

Let $g\in L^\infty$ be a piecewise constant centred global observable.
Using~\eqref{eq:p},
\[
\int_{Y_{j+i}}g\circ f^j\,d\mu=\sum_{p=1}^q\sum_{\psi(r)=p}\bar g_{i,p} \mu(Y_{j+i,r}).
\]
This combined with~\eqref{eq:K2} means that we have to show that $\lim_{n\to\infty} A_n=0$ where
\[
A_n=
\sum_{i\in R_n} \sum_{p=1}^q\bar g_{i,p}
\sum_{j\in S_{i,n}}
a_{n-j}^{-1}\sum_{\psi(r)=p}\mu(Y_{j+i,r}).
\]
Here, 
\[
R_n\subset\big\{\eps^2 n< i<(\eps^{-1}-1)n\big\}, \qquad
S_{i,n}\subset\big\{\eps \max\{i,n\}<j \le \min\{\eps^{-1}i, (1-\eps)n\}\big\}.
\]
(The containments might be strict since $i$ and $j$ are integers.)

\vspace{1ex} \noindent
\textbf{Step 1: Summation by parts over $j$.}
Let $\phi_{i,n}=\min S_{i,n}$, $\Phi_{i,n}=\max S_{i,n}$. 
Define
$\mu_{k,p}=\sum_{\psi(r)=p}\sum_{n\ge k}\mu(Y_{n,r})$.

Then
$A_n=
\sum_{i\in R_n}\sum_{p=1}^q\bar g_{i,p} B_{i,n,p}$ where
\[
B_{i,n,p}  = \sum_{j\in S_{i,n}} a_{n-j}^{-1}\sum_{\psi(r)=p}\mu(Y_{j+i,r})=
\sum_{\phi_{i,n}\le j\le \Phi_{i,n}} a_{n-j}^{-1}\{\mu_{j+i,p}-\mu_{j+i+1,p}\}.
\]
Summation by parts gives
\[
B_{i,n,p}  = 
\sum_{\phi_{i,n}\le j\le \Phi_{i,n}} \{a_{n-j}^{-1}-a_{n-j+1}^{-1}\}\mu_{j+i,p}
\;+\; \partial_{i,n,p,\phi_{i,n}} \;-\; \partial_{i,n,p,\Phi_{i,n}+1}
\]
where the boundary terms are 
\[
\partial_{i,n,p,c_{i,n}}=a_{n-c_{i,n}+1}^{-1}\,\mu_{c_{i,n}+i,p},
\qquad c_{i,n}=\phi_{i,n},\,\Phi_{i,n}+1.
\]

\vspace{1ex} \noindent
\textbf{Step 2: Simplification of $B_{i,n,p}$.}
Let $\tilde \mu_k=k^{-\alpha}$.
By~\eqref{eq:Y}, 
$\mu_{k,p}=\gamma_p\tilde\mu_k+o(k^{-\alpha})$.
Using also that $a_{n-c_{i,n}+1}^{-1}\ll n^{-(1-\alpha)}$,  
we obtain
\[
\partial_{i,n,p,c_{i,n}} = \gamma_p \tilde\partial_{i,n,c_{i,n}} +o(n^{-1}).
\qquad
\tilde\partial_{i,n,c_{i,n}} =a_{n-c_{i,n}+1}^{-1} \tilde\mu_{c_{i,n}+i}.
\]
Similarly, using in addition that
$a_{n-j}^{-1}-a_{n-j+1}^{-1} \ll
n^{-(2-\alpha)}$ and $\Phi_{i,n}-\phi_{i,n}\ll n$,
\[
\sum_{\phi_{i,n}\le j\le \Phi_{i,n}} \{a_{n-j}^{-1}-a_{n-j+1}^{-1}\}\mu_{j+i,p}
=\gamma_p \sum_{\phi_{i,n}\le j\le \Phi_{i,n}} \{a_{n-j}^{-1}-a_{n-j+1}^{-1}\}\tilde\mu_{j+i}
+o(n^{-1}).
\]
Hence, $B_{i,n,p}=\gamma_p \tilde B_{i,n}+o(n^{-1})$ where
\[
\tilde B_{i,n}  = 
\sum_{\phi_{i,n}\le j\le \Phi_{i,n}} \{a_{n-j}^{-1}-a_{n-j+1}^{-1}\}\tilde\mu_{j+i}
\;+\; \tilde \partial_{i,n,\phi_{i,n}} \;-\; \tilde\partial_{i,n,\Phi_{i,n}+1}
\]
Moreover, reversing the summation by parts argument in Step~1, we see immediately that
\[
\tilde B_{i,n}=
\sum_{\phi_{i,n}\le j\le \Phi_{i,n}} a_{n-j}^{-1}\{\tilde\mu_{j+i}-\tilde\mu_{j+i+1}\}=
\alpha \sum_{\phi_{i,n}\le j\le \Phi_{i,n}} a_{n-j}^{-1}(j+i)^{-(\alpha+1)}
+O(n^{-2}).
\]
Since 
$|R_n|\ll n$ and $|\bar g_{i,p}|\le |g|_\infty$, 
\begin{align*}
A_n & =\sum_{i\in R_n}\sum_{p=1}^q\bar g_{i,p}\Big\{\gamma_p\alpha \sum_{\phi_{i,n}\le j\le \Phi_{i,n}} a_{n-j}^{-1}(j+i)^{-(\alpha+1)} + o(n^{-1})\Big\}
\\ & =\alpha \sum_{i\in R_n}\bar g_i \sum_{\phi_{i,n}\le j\le \Phi_{i,n}} a_{n-j}^{-1}(j+i)^{-(\alpha+1)} \;+\; o(1).
\end{align*}

\vspace{1ex} \noindent
\textbf{Step 3: Summation by parts over $i$.}
We have reduced to showing that $\lim_{n\to\infty}C_n=0$ where
\[
C_n=\sum_{i\in R_n}\bar g_i \sum_{\phi_{i,n}\le j\le \Phi_{i,n}} a_{n-j}^{-1}(j+i)^{-(\alpha+1)} 
=\sum_{\eps n<j<(1-\eps)n} a_{n-j}^{-1}
\sum_{\eps j<i<\eps^{-1}j}\bar g_i (j+i)^{-(\alpha+1)} .
\]
Define
	\[
b_i=\sum_{k=n_0}^i k^{-\alpha}\bar g_k, \qquad  
b_n^*=\sup_{i\le n}|b_i|, \qquad
d_{i,j}= 
i^\alpha(j+i)^{-(\alpha+1)} .
\]
Then $d_{i,j}\ll n^{-1}$ and $d_{i,j}-d_{i+1,j}\ll n^{-2}$.
Also, 
$b_n^*=o(a_n)=o(n^{1-\alpha})$ by Lemma~\ref{lem:g}.

Summation by parts over $i$ gives
\begin{align*}
\sum_{\eps j<i<\eps^{-1}j} \bar g_i (j+i)^{-(\alpha+1)}
& =\sum_{\eps j<i<\eps^{-1}j} \bar g_i i^{-\alpha} d_{i,j}
=\sum_{\eps j<i<\eps^{-1}j} (b_i-b_{i-1}) d_{i,j}
\\ &  =\sum_{\eps j<i<\eps^{-1}j} b_i (d_{i,j}-d_{i+1,j})
\;+\; O(b_n^* n^{-1}) 
\\ &\ll n\cdot b_n^* \cdot n^{-2}+b_n^* n^{-1}=2 b_n^* n^{-1}=o(n^{-\alpha}).
\end{align*}
Hence $|C_n|\ll n\cdot n^{-(1-\alpha)}\cdot o(n^{-\alpha})=o(1)$,
completing the proof.
\end{proof}

\end{document}
